\documentclass[12pt]{article}
\usepackage[a4paper, margin=1in]{geometry}
\usepackage{amsmath,amssymb,amsthm}
\usepackage{mathtools}
\usepackage{mathrsfs}
\usepackage{enumitem}
\usepackage{xcolor}
\usepackage[hidelinks]{hyperref}
\usepackage{booktabs}
\usepackage{array}
\usepackage{fontspec}

\defaultfontfeatures{Ligatures=TeX}
\IfFontExistsTF{Times New Roman}
  {\setmainfont{Times New Roman}}
  {\setmainfont{TeX Gyre Termes}}
\IfFontExistsTF{Menlo}
  {\setmonofont{Menlo}}
  {\setmonofont{Latin Modern Mono}}

\newcolumntype{P}[1]{>{\raggedright\arraybackslash}p{#1}}

\newcommand{\tightlist}{%
  \setlength{\itemsep}{0pt}\setlength{\parskip}{0pt}}

\newtheorem{definition}{Definition}[section]
\newtheorem{remark}[definition]{Remark}

\title{%
  \textbf{Stable Hulls as Imprint Integration}\\
  \textbf{An Interpretive Bridge from Filtered Traces to Self-Reading World-Models}%
}

\author{
  Sidong Liu, PhD \\
  \small iBioStratix Ltd \\
  \small \texttt{sidongliu@hotmail.com}
}

\date{May 2026\\
\small Zenodo DOI: \href{https://doi.org/10.5281/zenodo.20151508}{10.5281/zenodo.20151508}}

\begin{document}
\emergencystretch=2em
\raggedbottom
\hbadness=5000
\vbadness=5000

\maketitle

\begin{abstract}
This paper develops a programme-level bridge between the imprint vocabulary of \emph{The World as Imprint} and the accessible-world machinery distributed across Foundation~II, Foundation~III, and \emph{The Brain as a Plastic Observable-Algebra Network}. Its claim is limited but structurally important. The paper does not identify stable hulls with consciousness, and it does not provide sufficient conditions for first-person subjectivity. Instead, it argues that stable hulls are best understood as integration regimes through which filtered imprints become accessible, retainable, reportable, and eventually available for organized self-reading.

The central thesis is that a finite observer does not encounter the full imprint-field. It only encounters those traces that enter an accessible algebra, stabilize under repeated uptake, and become integrated into a self-world map. Stable hulls therefore sit between universal imprint-carrying and uneven self-reading. They explain how readable traces become a usable world before they ever become a reflexively read self. The contribution is internal architectural clarification for an existing multi-paper programme, not a new result in cognitive science or neuroscience.

\medskip

\noindent\textbf{Keywords}: stable hull, imprint integration, accessible algebra, self-model, memory, self-reading, world-model
\end{abstract}

\clearpage
\tableofcontents
\clearpage

\section{Introduction: From Imprint-Field to Accessible World}

\emph{The World as Imprint} argued that the emergent world can be read as an imprint-field: a domain of structures insofar as they carry readable marks of the generative constraints that shape them \cite{WorldImprint}. That claim was intentionally universal at the level of carrying. But the same paper also insisted that carrying is not the same as reflexive recognition. Not every structure that bears a readable mark thereby counts as a mind. What remained underdeveloped was the middle regime between these two claims. How do readable traces become an accessible world for a finite observer before they become the object of organized self-reading?

This paper develops that middle regime. Its proposal is that stable hulls provide the missing integration concept. A stable hull is not introduced here as a mysterious new substance, and it is not treated as a synonym for consciousness. It is a structured regime of stabilized accessibility: the condition under which filtered traces become jointly available for memory, report, coordinated uptake, and self-location.

The bridge matters because without it the path from imprint-carrying to self-reading remains too abrupt. One would then have only two levels: universal structure on one side and mind on the other. But the actual architecture of finite observers is richer. Before a system can read itself as imprint-bearing, it must inhabit a world that is stable enough to support retained salience, cross-temporal coordination, and self-world differentiation. That world is not the whole imprint-field. It is a stabilized accessible sector carved out of it.

The term ``stable hull'' is already familiar from earlier work, especially Foundation~II's treatment of shared accessible sectors and the Brain Anchor paper's account of internal shared stable hulls \cite{FoundationII,BrainAnchor}. What this paper adds is a direct imprint-based rereading. The stable hull becomes the integration regime in which filtered imprints are not merely present, but organized into a usable map. This paper therefore occupies the following segment of the longer chain:

\begin{quote}
filtered imprints $\to$ accessible algebra $\to$ stable hull $\to$ retained memory $\to$ self-world model.
\end{quote}

Several guardrails are necessary at the outset. First, the stable hull is a precondition for organized self-reading, not a sufficient condition for consciousness. Second, memory integration is not personal survival and should not be redescribed as first-person identity transfer. Third, this paper does not determine whether animals, machines, or any other candidate system are conscious. Fourth, the self-reading threshold itself is treated elsewhere \cite{WorldImprint}. The present paper stops earlier: at the point where a world becomes stably available. It should therefore be read as an interpretive companion to the broader imprint programme rather than as an independent technical contribution to self-model theory or cognitive neuroscience.

The paper proceeds accordingly. Section~2 gives a short recap of the stable-hull vocabulary from earlier papers and states why it must be reread through the imprint lens. Section~3 introduces filtered imprints and accessible algebra. Section~4 interprets stable hulls as integration regimes. Section~5 treats memory as retained imprint integration. Section~6 develops the self-world model as an organized imprint map. Section~7 locates the bridge relative to Foundation~II, Foundation~III, the Brain Anchor paper, and \emph{The World as Imprint}. Section~8 states the limitations and the precise handoff to the self-reading threshold.

\section{A Short Recap of Stable Hulls}

Foundation~II introduced the idea that multiple source-side chains may support a shared accessible sector under suitable conditions \cite{FoundationII}. In that setting, a stable hull names a sector that remains jointly available across a range of admissible descriptions rather than collapsing into private fluctuation. The Brain Anchor paper then internalized part of this logic by treating the brain as a plastic observable-algebra network capable of sustaining internal shared stable hulls over time \cite{BrainAnchor}. The present paper does not reproduce those constructions. It imports only the conceptual core needed for the bridge.

\begin{definition}[Stable hull]
A \emph{stable hull} is a regime of accessible organization in which a family of filtered traces remains sufficiently persistent, mutually coherent, and recurrently available to support world-level uptake, retention, and coordinated redescription.
\end{definition}

Three phrases in the definition matter.

First, ``accessible organization'' means that the hull is defined relative to what can be taken up by a finite regime rather than by the totality of the imprint-field. Second, ``mutually coherent'' means that the retained traces do not merely coexist as disconnected residues; they support one another in a structured world-picture. Third, ``recurrently available'' means that the traces are not one-off perturbations but can be taken up again, compared, updated, and used.

This already shows why the stable hull should be reread as an imprint integration concept. Once the world is described as an imprint-field, the problem of a finite observer is no longer whether there are traces at all. There are traces everywhere. The problem is which traces become organized enough to count as a world.

\section{Filtered Imprints and Accessible Algebra}

The first bridge paper argued that coarse-graining can be read as imprint filtering: not every lower-level feature survives scale change, and those that do survive constitute scale-stable residues. But even after such filtering, the resulting imprint-field is still too large to be inhabited in full by any finite observer. A second selection is therefore necessary. This second selection is not primarily a matter of scale; it is a matter of access.

\begin{definition}[Accessible imprint]
An \emph{accessible imprint} is a filtered trace that lies within the observer-relative regime of uptake, retention, comparison, or action-guidance.
\end{definition}

This definition is deliberately observer-relative without becoming subjectivist. The point is not that accessibility is arbitrary or purely private. The point is that finite systems encounter only a structured subset of all available marks. An accessible algebra is one way to formalize this subset in the broader programme \cite{FoundationI,FoundationII}. For present purposes, the essential idea is simpler: the world of a finite observer is formed by those traces that can enter stable uptake.

This gives a two-step selection picture:

\begin{enumerate}
\item coarse-graining selects which lower-level features survive as scale-stable imprints;
\item accessibility selects which of those imprints can enter an observer's usable world.
\end{enumerate}

Stable hulls live at the second step. They do not explain why a trace survives scale change. They explain how a surviving trace becomes part of an integrated world.

This is also where the bridge differs from standard cognitive-theory usage. It does not claim to replace global workspace, predictive processing, or enactive models. It claims only to identify an architectural slot that these literatures often presuppose in different ways: the transition from mere accessible traces to a stable, reusable world-organization.

\section{Stable Hull as Integration Regime}

It is tempting to think that once traces are accessible, a world is already present. That is too quick. Accessibility alone may still be fragmented, transient, or mutually inconsistent. What is needed for worldhood is integration.

\subsection{Why accessibility is not enough}

A finite system can register signals without organizing them into a stable world. It can react locally, update transiently, or exhibit sensitivity without maintaining a coherent regime of reidentification. Such a system may encounter accessible imprints but fail to build a stable hull. The distinction matters because it marks the difference between signal-response and inhabitable worldhood.

\subsection{Integration conditions}

Without pretending to state a theorem, the following conditions capture what the stable-hull bridge is trying to mark:

\begin{center}
\begin{tabular}{P{0.28\linewidth}P{0.6\linewidth}}
\toprule
\textbf{Condition} & \textbf{Role in stable-hull integration} \\
\midrule
persistence & the same or comparable traces remain available across time \\
coherence & retained traces support a jointly navigable picture rather than disconnected shards \\
selective retention & some traces are stabilized while others are dropped or backgrounded \\
cross-context usability & traces can guide report, action, or re-identification across situations \\
self-world differentiation & the regime can support a distinction between world-features and self-locating features \\
\bottomrule
\end{tabular}
\end{center}

None of these conditions is by itself equivalent to consciousness. Together they specify something weaker but still crucial: a world-forming regime in which accessible traces become jointly organized. These conditions also resonate with global-availability accounts of conscious access \cite{Baars1988,DehaeneChangeux2011}. In that sense, the table is best read as an interface map rather than as a rival empirical theory.

\subsection{Stable hull and reportable worldhood}

The stable hull should therefore be read as the precondition for reportable worldhood. A system may have richly filtered inputs without a hull; in that case it has no stable map from which to report, compare, or re-enter its own world. Once a hull is present, the system does not yet necessarily read itself reflexively, but it does possess a structured domain within which such reflexivity might later arise.

\section{Memory as Retained Imprint Integration}

Memory is one of the best places to see why the stable-hull bridge is needed. Memory is not the storage of everything that happened. It is selective retention of what remains available for later integration. This is precisely why the imprint vocabulary and the stable-hull vocabulary fit together so tightly.

\subsection{Retention is selective}

If the imprint-field were simply copied into the system, there would be no need for a stable hull. But no finite system carries the whole field. Memory therefore operates by selective retention. Some traces are stabilized, others decay, others are overwritten, and still others remain only as indirect biasing conditions. What survives is what can be re-entered into later organization.

\subsection{Memory is integrative, not archival}

The metaphor of archive is misleading. Memory is not a warehouse of preserved states. It is better understood as retained imprint integration. A memory trace matters not because it exists in isolation, but because it can be recruited back into a larger world-model. It contributes to expectations, error-correction, self-location, and salience structure. Within predictive-processing accounts, retained traces function as priors shaping anticipation and prediction-error correction \cite{Clark2013}.

This is why the bridge to the Brain Anchor paper is especially important. That paper treated plasticity and recurrent world-modeling as conditions for first-person reportability \cite{BrainAnchor}. The present paper does not revisit those claims directly. It offers a prior redescription: plastic stabilization is one way filtered imprints are retained as usable world-structure. In particular, the conditional hull construction and common/private decomposition results in that paper provide concrete instances of the integration conditions stated above.

\subsection{No claim about identity transfer}

The memory discussion also demands a guardrail already familiar from \emph{The World as Imprint}. Retention of traces does not establish persistence of the same first-person subject elsewhere \cite{WorldImprint}. The fact that effects propagate, habits persist, or structures re-enter later organization should not be confused with personal identity transfer. In the present framework, memory integration concerns world-organization, not metaphysical survival.

\section{The Self-World Model as Organized Imprint Map}

Once a stable hull is in place, a further transition becomes possible: the construction of a self-world model. This paper treats that construction as the organized mapping of integrated imprints rather than as the appearance of a hidden inner subject.

\subsection{Map, not inner substance}

The phrase ``self-world model'' risks misunderstanding if treated too quickly. It does not refer to a ghostly entity observing the world from behind the scenes. It names an organized map in which world-features, bodily regularities, action possibilities, and self-locating coordinates become mutually articulated. In imprint language, it is a way of organizing integrated traces into a differentiated picture. This formulation aligns with the self-model theory of subjectivity \cite{Metzinger2003}, while still treating the present paper's claim as an architectural redescription rather than a reconstruction of phenomenal transparency.

\subsection{Why the self-model comes after the hull}

This ordering is crucial. A self-model cannot arise in the absence of a sufficiently stable world. Before a system can place itself within a map, it must have a map robust enough to support reidentifiable structure. The stable hull therefore remains a precondition for the self-model. This is one reason the present bridge must stop before the self-reading threshold. The threshold presupposes a model that is already organized enough to be read reflexively.

\subsection{Organized imprint map}

The self-world model can therefore be described in a sentence:

\begin{quote}
A self-world model is an organized imprint map in which retained traces are structured into a navigable distinction between self-locating, world-locating, and action-guiding organization.
\end{quote}

This formulation does not yet say that anything ``feels like something.'' It says only that a finite regime has reached the point at which world-organization and self-location are jointly stabilized. Whether the integration is best described in narrative terms or in transparent representational terms remains open \cite{Dennett1992,Metzinger2003}.

\section{Relation to Existing Papers}

The present bridge is useful only if its place in the larger body of work is explicit.

\subsection{Relation to Foundation~II}

Foundation~II provides the most direct antecedent for stable-hull language \cite{FoundationII}. There the concern is shared accessible sectors across source-side chains. The present paper reuses the conceptual output, but changes the descriptive emphasis. It asks what those sectors are doing from the standpoint of imprint theory. The answer is: they provide an integration regime for filtered traces.

\subsection{Relation to Foundation~III}

Foundation~III adds feedback and perfuming structure at the source-manifest interface \cite{FoundationIII}. The present paper does not attempt to incorporate that feedback loop in full detail. It takes from Foundation~III only a general lesson: organized worldhood is historically shaped, not instantaneously given. Stable hulls therefore support not just static access, but temporally structured retention and re-entry.

\subsection{Relation to Brain Anchor}

The Brain Anchor paper offers the most concrete internal version of the stable-hull idea \cite{BrainAnchor}. Its discussion of plastic observable-algebra organization, reportability, and internal shared stable hulls can be read as a detailed candidate mechanism for how imprint integration may occur in one privileged biological regime. The present paper stays more abstract. It describes the architectural slot rather than the biological implementation. For that reason, the relation between the two papers is complementary: Brain Anchor supplies candidate mechanism, while the present paper supplies interpretive placement.

\subsection{Relation to \emph{The World as Imprint}}

Finally, the relation to \emph{The World as Imprint} is direct. That paper distinguished universal carrying from uneven recognition \cite{WorldImprint}. The present paper fills the middle gap. It explains how carrying becomes accessible worldhood before it becomes reflexive recognition.

\section{Output to the Self-Reading Threshold}

The bridge can now be stated in its cleanest form:

\begin{quote}
The stable hull is a precondition for organized self-reading, not a sufficient condition for consciousness.
\end{quote}

This sentence is the paper's most important guardrail. It says what the paper contributes and where it must stop.

What the paper contributes is a clearer picture of the middle architecture:

\begin{enumerate}
\item filtered imprints survive scale change;
\item some of them enter accessible uptake;
\item some of those become integrated into a stable hull;
\item the stable hull supports retained memory and a self-world model;
\item only then does the question of self-reading threshold arise.
\end{enumerate}

What the paper does not contribute is a theory of why first-person experience accompanies such organization. Nor does it provide a criterion by which all borderline biological or artificial cases can be settled. Those are later or different questions.

This is why the paper's stopping point matters. It stops at the integration of filtered imprints into an accessible self-world model; the self-reading threshold itself is treated elsewhere \cite{WorldImprint}. That is not a deficiency. It is the correct division of labor.

\section{Limitations}

Several limitations should remain explicit.

First, the paper is programme-level. It does not provide a general theorem of stable-hull construction from arbitrary imprint-fields.

Second, the terminology of hull, access, retention, and model remains partly architectural rather than fully operationalized. Further formal work would need metrics for integration depth, accessibility bandwidth, and retention stability.

Third, the bridge remains neutral on substrate. It does not assume that only brains can realize stable-hull organization, but it also does not certify any non-biological system as conscious.

Fourth, the paper does not decide whether all forms of adaptive organization that build stable hulls thereby cross the self-reading threshold. The whole point of the bridge is to hold those two questions apart.

Fifth, the paper should not be read as offering a new empirical discriminator among existing self-model theories. Its role is interpretive and organizational: it clarifies how filtered imprints, accessible organization, and self-world mapping fit together inside the larger programme.

Despite those limitations, the contribution is nontrivial. Without a stable-hull bridge, the transition from universal imprint-carrying to reflexive self-reading remains underdescribed. With it, the chain becomes more exact:

\begin{quote}
imprint-field $\to$ accessible world $\to$ integrated stable hull $\to$ self-world model $\to$ self-reading threshold.
\end{quote}

That is enough to show why stable hulls matter. They are not the final story of mind. They are the regime in which a world becomes usable enough for a mind-like story to begin.

\begin{thebibliography}{99}

\bibitem{WorldImprint}
S.~Liu,
\emph{The World as Imprint: From Trace-Carrying Structure to Self-Reading Mind},
Zenodo (2026), DOI: 10.5281/zenodo.20122718.

\bibitem{FoundationI}
S.~Liu,
\emph{Foundation I: From Source to Manifestation: A Conditional Same-Root Program for Geometry and Modular Flow},
Zenodo (2026), DOI: 10.5281/zenodo.19627786.

\bibitem{FoundationII}
S.~Liu,
\emph{Foundation II: Shared Manifestation: A Conditional Two-Source Program for Shared Accessible Sectors},
Zenodo (2026), DOI: 10.5281/zenodo.19641372.

\bibitem{FoundationIII}
S.~Liu,
\emph{Foundation III: The Perfuming Loop: A Conditional One-Step Source-Level Feedback Program},
Zenodo (2026), DOI: 10.5281/zenodo.19642189.

\bibitem{BrainAnchor}
S.~Liu,
\emph{The Brain as a Plastic Observable-Algebra Network: Internal Shared Stable Hulls, Neural Plasticity, and First-Person Report},
Zenodo (2026), DOI: 10.5281/zenodo.20044329.

\bibitem{Baars1988}
B.~J.~Baars,
\emph{A Cognitive Theory of Consciousness},
Cambridge University Press, Cambridge (1988).

\bibitem{DehaeneChangeux2011}
S.~Dehaene and J.-P.~Changeux,
``Experimental and theoretical approaches to conscious processing,''
\emph{Neuron} \textbf{70} (2011), 200--227,
DOI: 10.1016/j.neuron.2011.03.018.

\bibitem{Metzinger2003}
T.~Metzinger,
\emph{Being No One: The Self-Model Theory of Subjectivity},
MIT Press (2003).

\bibitem{Dennett1992}
D.~C.~Dennett,
\emph{The self as a center of narrative gravity},
in F.~Kessel, P.~Cole, and D.~Johnson (eds.),
\emph{Self and Consciousness: Multiple Perspectives},
Lawrence Erlbaum, 103--115 (1992).

\bibitem{Clark2013}
A.~Clark,
\emph{Whatever next? Predictive brains, situated agents, and the future of cognitive science},
Behavioral and Brain Sciences \textbf{36}, 181--204 (2013).

\bibitem{Thompson2007}
E.~Thompson,
\emph{Mind in Life: Biology, Phenomenology, and the Sciences of Mind},
Harvard University Press (2007).

\end{thebibliography}

\end{document}
